\chapter{Our Constitution}
\section{Get Rolling}
\textbf{Write two rules that you follow:}
\vspace{0.1in}
\par
At School:
\begin{enumerate}
    \item \textbf{Don not push and run} in the corridors and staircases
    \item \textbf{Listen to your teacher} and follow classroom instructions
\end{enumerate}
In the playground:
\begin{enumerate}
    \item \textbf{Wait for your turn} on the swings and slides
    \item \textbf{Do not push or shove} other children while playing
\end{enumerate}
On road:
\begin{enumerate}
    \item \textbf{Always walk on the footpath} (pavement).
    \item \textbf{Cross the road at the zebra crossing} after looking right and left.
\end{enumerate}
\section{What to Know}
A constitution is a set of rules that guides everyone in a country. It grants rights, duties and liabilities to the citizens of a country.
\section{\textbf{Learning Objectives}}
\begin{itemize}
    \item to explain the meaning of the rule of law
    \item to define a constitution and discuss the importance
    \item to describe the making of the Indian Constitution
    \item to analyse the ideas given in the Preamble to the Indian Constitution
    \item to discuss the features of a sovereign, socialist, secular and democratic republic
    \item to list the fundamental rights and duties of Indian citizens
\end{itemize}
\section{I Explore}
\section{Rules or Laws}
Rules provide guidelines to people. It is important to make and follow rules in games, organisations, schools and countries. You may have read the rules to be followed at school in your school almanac.
\vspace{0.1in}
\par
Similarly, most countries have a set of rules that are to be followed by all. It is important to have rules and regulations for the safety and efficient functioning of a country. Everyone has to abide by these rules. This set of rules is called the \textbf{constitution} of a country and it is accepted by all people living together in the country. A constitution is a set of laws that describe the rights and duties of citizens. Our country also has a constitution.
\vspace{0.1in}
\par
Law is a supreme and no human being has authority above it. The \textbf{rule of law} is guaranteed by the constitution, which means that all authorities follow it strictly. All citizens of India are governed by the rule of law. It means that all laws apply equally to all citizens of the country. No one is above the law.
\section{Making of the Indian Constitution}
The Constitution of India is the supreme law of the nation. It clearly states how the Government of India would be formed, what work the government would do and what powers it would have. The Constitution also states the rights and duties of the citizens of India.
\vspace{0.1in}
\par
The constituent Assembly was formed on 9 December 1946 to draft the Constitution of India. Dr Rajendra Prasad was the President of the Constituent Assembly.
\vspace{0.1in}
\par
On 29 August 1947, the Constituent Assembly set up a Drafting Committee under the chairmanship of Dr BR Ambedkar. The objective of the drafting committee was to prepare a draft constitution of India.
\vspace{0.1in}
\par
Most members of this committee were involved in the freedom struggle against British rule. They wanted the Constitution of India to reflect the hopes of Indians, who after gaining freedom from the foreign rulers, wanted a goverment that would treat them as equals in social, cultural and political fields.
\vspace{0.1in}
\par
The Constitution of India came into effect on 26 January 1950 and India became a republic. Every year, we celebrate this day as the Republic Day.
\vspace{0.6in}
\section{Preamble to the Indian Constitution}
The Constitution of India opens with the \textbf{Preamble}. The word 'preamble' means introduction. The Preamble to the Indian Constitution is a brief introduction of the fundamental values and the philosophy on which the Indian Constitution is based. It is considered a key to the Constitution, as it reflects the objectives and the ideals of the Constitution.
\subsection{Values Given in the Preamble to the Indian Constitution}
The Preamble to the Indian Constitution starts with \textbf{WE, THE PEOPLE OF INDIA, having solemnly resolved to constitute India into a SOVEREIGN, SOCIALIST, SECULAR, DEMOCRATIC REPUBLIC...}
\vspace{0.1in}
\par
'We the people of India' means that the people of India have full right to make decisions of all matters related to India.
\vspace{0.1in}
\par
\textbf{'SOVEREIGN'} means the India is a free country. No country or individual can challenge India's independence from other nations and external powers. It also means that the people of India have full right to make decision on all matters related to India.
\vspace{0.1in}
\par
In 1976, the Constitution was amended, or changed, and the two terms \textbf{socialist} and \textbf{secular} were added to the Preamble. India became a \textbf{'SOCIALIST'} country, meaning that everyone in society should have equal access to the wealth produced in India. Our constitution helps the attainment of the ideals of:
\begin{itemize}
    \item removal of inequalities from the society
    \item basic necessities to all
    \item equal pay for equal work
\end{itemize}
On Constitution declares that India is a \textbf{'SECULAR'} country, meaning that all the religions in India get equal respect and protection from the State. Our country is not guided by any religion. There is no State or official religion. Everyone is free to believe in and follow any religion. The government will not discriminate against any individual or group on the basis of religion.
\vspace{0.4in}
\par
The term \textbf{'DEMOCRATIC'} implies that all the adult citizens of India have the right to elect their representatives. The representatives rule the country and are responsible to the people of India. The democratic principles incorporated in the Constitution are universal adult franchise (popularly known as one person one vote), elections at regular intervals, Fundamental Rights, etc.
\vspace{0.1in}
\par
As opposed to a monarchy, the Indian Constitution defines our nation as a republic. The term \textbf{'REPUBLIC'} indicates that the head of the State, or the president of India, is elected by the people indirectly for a fixed period. He or she is not a hereditary ruler like a king or a queen.
\vspace{0.3in}
\par
I wonder how the principles of 'sovereign, socialist, secular and democratic republic' outlined in the Preamble to the Indian Constitution have shaped the values and functioning of our country!
\subsection{Objectives of the Indian Constitution}
The Preamble lists four main aims which are to be achieved by India for all its citizens:
\begin{itemize}
    \item JUSTICE, social, economic and political
    \item LIBERTY of thought, expression, belief, faith and worship
    \item EQUALITY of status and of opportunity
    \item FRATERNITY, assuring the dignity of an individual and the unity and intergrity of the nation
\end{itemize}
The constitution is the supreme law and it helps:
\begin{itemize}
    \item to maintain integrity in society
    \item to promote unity among citizens to build a great nation
\end{itemize}
The main objective of the Indian Constitution is to promote harmony throughout the nation. This objective can be achieved with the help of the following means:
\vspace{0.1in}
\par
\textbf{Justice:} The citizens of India cannot be treated unequally because of differences in their caste, religion, gender, income and wealth.
\vspace{0.1in}
\par
\textbf{Liberty:} The citizens of India are free to choose their way of life. They have freedom on thought, expression, belief, faith and worship.
\vspace{0.1in}
\par
\textbf{Equality:} Every citizen of India is entitled to equality before law. No section of society has any special privileges and all the people are given equal opportunities for everything without any discrimination.
\vspace{0.1in}
\par
\textbf{Fraternity:} The term 'fraternity' means a feeling of brotherhood and an emotional attachment with the country and all its people. Fraternity helps to promote dignity and unity in the nation.
\vspace{1in}
\section{Our Fundamental Rights}
The Constitution of India guarantees citizen of the country some basic rights. These rights are called \textbf{Fundamental Rights.} These rights ensure the well-being of the citizens. If any of these rights are denied to us anywhere in India, we can move to court for their protection. There are six Fundamental Rights granted to the citizens of India.
\vspace{1.1in}
\textbf{Right to Equality:} The rights to equality means that all Indian citizens are equal before law. All citizens irrespective of their religion, caste, gender, place of birth, etc. are treated equally by law.
\vspace{0.1in}
\par
\textbf{Right to Freedom:} The right to freedom allows all citizens the freedom:
\begin{itemize}
    \item of speech and expression
    \item to pursue any profession or occupation
    \item to form an assembly or association
    \item to live and travel anywhere within the country
    \item to do any kind of work at any place of their choice.
    \item of free and compulsory education for all children between 6 and 14 years of age.
\end{itemize}
\textbf{Right against Exploitation:} The right against exploitation ensures that nobody is forced to work against exploitation ensures that nobody is forced to work against his/her wish. The right does not allow any child below 14 years of age to be employed.
\vspace{0.1in}
\par
\textbf{Right to Freedom of Religion:} The right to freedom o religion means that all people are free to practice any religion that they choose.
\vspace{0.1in}
\par
\textbf{Cultural and Educational Rights:} This fundamental right protects minorities. The minorities can establish educational institutions of their choice to protect their language and culture.
\vspace{0.1in}
\par
\textbf{Right to Constitutional Remedies:} The right to Constitutional Remedies allows every citizen to go to court if any of their fundamental rights are talken away.
\par
In times of war or their emergencies, the government can limit or cancel some of these rights for some time.
\vspace{1.3in}
\section{Our Fundamental Duties}
The Indian Constitution mentions a list of duties that all citizens of India must fulfil. Such duties are called \textbf{Fundamental Duties.} These duties inspire the citizens to live in harmony and respect society. Some of the Fundamental Duties are as follows.
\begin{itemize}
    \item We must abide by the Constitution.
    \item We must respect our national flag and national anthem.
    \vspace{0.4in}
    \item We must protect the sovereignty and integrity of India.
    \item We must defend the country and render national services when called upon.
    \item We must safeguard public property and avoid violence.
    \item We must try to protect our culture.
    \item We must protect our natural environment.
    \item Parents or guardians must send children in the age group of 6-14 years to school.
\end{itemize}
The Constitution of India outlines the powers and limitations of the government as well as the rights and duties of citizens. It sets the framework for the governance of the country. It is a very important document, as it safeguards the country from dictatorship or abuse of power.
\vspace{0.3in}
\par
I wonder how the Fundamental Duties outlined in the Constitution of India contribute to fostering harmony and a sense of responsibility among citizens!
\vspace{2.2in}
\section{Now I Know}
\begin{itemize}
    \item Rules and laws provide guidelines to people.
    \item A constitution is a set of laws that describe the rights and duties of citizens.
    \item The Constitution of India is the supreme law of the nation. It was framed by a representative body known as the Cpnstituent Assembly.
    \item The Indian Constitution came into effect on 26 January 1950. From that day, our country came to be known as the Republic of India. Every year, we celebrate 26 January as the Republic Day of India.
    \item The Preamble describes India as a SOVEREIGN, SOCIALIST, SECULAR, DEMOCRATIC, REPUBLIC.
    \item The Preamble lists four main aims which India should achieve for all its citizens. These are justice, liberty, equality, and fraternity.
    \item Our Fundamental Rights are the basic rights. They are guaranteed and protected by the Constitution of India.
    \item The Indian Constitution mentions a list of duties that all citizens of India must fufil. These are called Fundamental Duties.
    \item The Constitution of India outlines the powers and limitations of the government as well as the rights and duties of the citizens.
\end{itemize}